\documentclass[10pt]{article}
\usepackage[T1]{fontenc}
\usepackage[utf8]{inputenc}
\usepackage[margin=2.3cm]{geometry}
\usepackage{amsmath,booktabs,hyperref}
\hypersetup{colorlinks=true,linkcolor=black,urlcolor=blue,citecolor=black}
\title{A modular approach to selective molecular sensing}
\author{Priya Nair\thanks{p.nair@example.org} \and Thomas Green}
\date{}
\begin{document}
\maketitle
\begin{abstract}
We describe a modular sensing workflow in which calibration, selectivity, and
repeatability are reported as separate experimental questions. The starter
shows the compact narrative, data table, and reference structure expected of
an RSC-style manuscript.
\end{abstract}

\section{Introduction}
Selective detection depends on both the response of the active material and
the conditions under which the response is measured. Reporting those factors
separately helps readers compare a new system with established methods.

\section{Experimental}
Solutions were prepared from analytical-grade reagents and diluted with
deionized water. The response ratio was calculated as
\begin{equation}
  R=\frac{I-I_0}{I_0},
  \label{eq:response}
\end{equation}
where $I$ is the measured intensity and $I_0$ is the blank response. Include
instrument settings, purification details, and safety information here.

\section{Results and discussion}
The reporting pattern in Table~\ref{tab:calibration} separates calibration
from selectivity. Replace the illustrative values with the measurements from
the completed experiment and provide a figure file when needed.
\begin{table}[ht]
  \centering
  \caption{Illustrative calibration summary.}
  \label{tab:calibration}
  \begin{tabular}{lrr}
    \toprule
    Analyte & Range (\textmu M) & $R$ at midpoint \\
    \midrule
    A & 0.5--10 & 0.42 \\
    B & 1--20 & 0.08 \\
    C & 1--20 & 0.05 \\
    \bottomrule
  \end{tabular}
\end{table}
The response in equation~\eqref{eq:response} should be accompanied by an
estimate of uncertainty and the number of independent replicates.

\section{Conclusions}
Separating calibration and selectivity makes the experimental claim easier to
evaluate. Use the journal-specific communication template when the manuscript
type requires it, and follow the current RSC author guidance.

\section*{Conflicts of interest}
There are no conflicts of interest to declare.

\begin{thebibliography}{9}
\bibitem{rsc} Royal Society of Chemistry, ``Article templates,'' publisher
author guidance, accessed 2026.
\bibitem{principles} D. A. Skoog, F. J. Holler and S. R. Crouch,
\emph{Principles of Instrumental Analysis}, Cengage Learning, 2018.
\end{thebibliography}
\end{document}
